% Source: Weinberg GR, physical PDF pp. 385--386
%         (printed pp. 364--365).
% Coverage: Section 12.4, The Gravitational Action; equations
%           (12.4.1)--(12.4.3).
% Figures/tables/footnotes: none.
% Status: source-reviewed and compile-clean.
% Uncertainties: the curvature/action sign conversion in (12.4.2)--(12.4.3)
%                is documented below and checked against Einstein's equation.
%
% MODERNIZATION CHECK: At fixed slots the target Riemann, Ricci, scalar, and
% Einstein tensors are the negatives of Weinberg's. Consequently the source
% minus sign in (12.4.2) becomes a plus, while its covariant-metric variation
% (12.4.3) becomes a minus. The resulting equation is G^{mu nu}=8 pi G
% T^{mu nu}, as required by the binding convention.
% LITERAL-OPERATOR AUDIT: both terms in the Palatini divergence, the minus
% sign in (12.4.3), and the matter term in the final field equation were
% checked independently against physical PDF p. 385 after conversion.

\subsection{The Gravitational Action}\label{sec:12.4}

So far, in this chapter the gravitational field \(g_{\mu\nu}\) has been an
external field that could be prescribed at will. (Indeed,
Eq.~\eqref{eq:12.2.2} usually provides the most convenient definition of the
energy-momentum tensor even in the absence of gravitation.) We will now give
\(g_{\mu\nu}\) field equations of its own by adding to the total action \(I\)
a purely gravitational term \(I_G\):
\begin{equation}
  I=I_M+I_G,
  \tag{12.4.1}\label{eq:12.4.1}
\end{equation}
\begin{equation}
  I_G
  =
  \frac1{16\pi G}\int\dd^4x\,\sqrt{-g(x)}\,R(x).
  \tag{12.4.2}\label{eq:12.4.2}
\end{equation}
Clearly \(I_G\) is a scalar, so this would be a good candidate for a theory
of gravitation even if we had no experience with gravitational phenomena. We
shall now show that the application to \(I\) of the principle of stationary
action does in fact yield Einstein's theory.

The curvature scalar \(R\) is defined as \(g^{\mu\nu}R_{\mu\nu}\), so a
variation \(\delta g_{\mu\nu}\) in the metric produces in the integrand of
Eq.~\eqref{eq:12.4.2} a change
\[
  \delta\!\left(\sqrt{-g}\,R\right)
  =
  \sqrt{-g}\,R_{\mu\nu}\delta g^{\mu\nu}
  +R\,\delta\sqrt{-g}
  +\sqrt{-g}\,g^{\mu\nu}\delta R_{\mu\nu}.
\]
According to Eq.~\eqref{eq:10.9.2}, converted to the curvature convention of
this edition, the change in the Ricci tensor is
\[
  \delta R_{\mu\nu}
  =
  \nabla_\lambda\delta\Gamma^\lambda{}_{\nu\mu}
  -\nabla_\nu\delta\Gamma^\lambda{}_{\lambda\mu},
\]
the covariant derivatives being defined as if
\(\delta\Gamma^\lambda{}_{\mu\nu}\) were a tensor (as in fact it is). Thus the
last term in \(\delta(\sqrt{-g}\,R)\) is
\[
\begin{aligned}
  \sqrt{-g}\,g^{\mu\nu}\delta R_{\mu\nu}
  =\sqrt{-g}\big[
    &\nabla_\lambda(g^{\mu\nu}\delta\Gamma^\lambda{}_{\nu\mu})
    \\
    &-\nabla_\nu(g^{\mu\nu}\delta\Gamma^\lambda{}_{\lambda\mu})
  \big],
\end{aligned}
\]
or, using Eq.~\eqref{eq:4.7.7},
\[
\begin{aligned}
  \sqrt{-g}\,g^{\mu\nu}\delta R_{\mu\nu}
  ={}&
  \partial_\lambda\!\left(
    \sqrt{-g}\,g^{\mu\nu}\delta\Gamma^\lambda{}_{\nu\mu}
  \right)
  \\
  &-\partial_\nu\!\left(
    \sqrt{-g}\,g^{\mu\nu}\delta\Gamma^\lambda{}_{\lambda\mu}
  \right).
\end{aligned}
\]
This term therefore drops out when we integrate over all space. Also,
\[
  \delta\sqrt{-g}
  =
  \frac12\sqrt{-g}\,g^{\mu\nu}\delta g_{\mu\nu},
  \qquad
  \delta g^{\mu\nu}
  =
  -g^{\mu\rho}g^{\nu\sigma}\delta g_{\rho\sigma},
\]
so the change in the action \eqref{eq:12.4.2} is
\begin{equation}
  \delta I_G
  =
  -\frac1{16\pi G}\int\dd^4x\,\sqrt{-g}\,
  \left(R^{\mu\nu}-\frac12g^{\mu\nu}R\right)
  \delta g_{\mu\nu}.
  \tag{12.4.3}\label{eq:12.4.3}
\end{equation}
Combining Eq.~\eqref{eq:12.4.3} with Eq.~\eqref{eq:12.2.2}, we see that the
total action \(I\) is stationary with respect to arbitrary variations in
\(g_{\mu\nu}\) if and only if
\[
  R^{\mu\nu}-\frac12g^{\mu\nu}R-8\pi G\,T^{\mu\nu}=0,
\]
which, of course, is the Einstein field equation.

[As another application of Eq.~\eqref{eq:12.4.3}, we may use it to derive the
contracted Bianchi identity. Since \(I_G\) is a scalar, it must be stationary
with respect to the variation \eqref{eq:12.3.1} in \(g_{\mu\nu}\). Repeating
the reasoning that led before to Eq.~\eqref{eq:12.3.2}, we now find that
\[
  \nabla_\nu
  \left(R^\nu{}_\lambda-\frac12\delta^\nu{}_\lambda R\right)=0,
\]
which we recognize as the contracted Bianchi identity \eqref{eq:6.8.3}.]

This formalism suggests that Einstein's theory might be modified by adding to
\(R\) in Eq.~\eqref{eq:12.4.2} terms proportional to \(R^2\), \(R^3\), etc.
As discussed in Section~\ref{sec:7.1}, such terms would only show up on a
sufficiently small space-time scale.
